% ======================================================================
% Порядок выполнения. Ответственный: A (собирает строки от всех).
% Требования: main.tex, строки 1851–1853 — краткая последовательность
% фактически выполненных действий разделов 1–3; какие пункты со звёздочкой выполнены.
% Образец: report_template.tex, строки 78–88.
% Каждый владелец присылает A одну строку про свою часть (или правит свою строку через PR).
% ======================================================================
\taskhead{Порядок выполнения}
\begin{enumerate}
\item[0] \todo{A: подготовка окружения (Python, библиотеки, XeLaTeX).}
\item[1.1] \todo{D: что сделано — аналитика, обе квадратуры, таблицы, графики.}
\item[1.2] \todo{D: что сделано — аналитика, замена переменных, якобиан, численный расчёт.}
\item[2.1] \todo{D: что сделано.}
\item[2.2] \todo{D: что сделано.}
\item[2.3] \todo{D: что сделано.}
\item[3.1] \todo{A: что сделано.}
\item[3.2] \todo{B: что сделано.}
\item[3.3] \todo{B: что сделано.}
\item[3.4] \todo{A: что сделано.}
\item[3.5] \todo{A/B/C: численные эксперименты 3.5.1–3.5.5.}
\item[3.6] \todo{C: что сделано.}
\end{enumerate}
\todo{A: указать, какие пункты со звёздочкой (3.1.6*, 3.2.5*) выполнены, или что они не выполнялись.}
